\section{The Multi-Component, Physics-Constrained Loss Function}

The training of the PINN is guided by a composite loss function that simultaneously optimizes for three distinct yet complementary objectives:
fidelity to observed data, adherence to the governing laws of physics, and conservation of total system energy.
This multi-component objective function,
as detailed in the Augmented Lagrangian Method (ALM) with a separate energy regularization term, enables the network to learn solutions that are not only accurate
but also physically plausible and generalizable beyond the specific training data points.
Specifically, the total loss function combines the data fidelity term $\mathcal{L}_{\text{data}}$ and the physics adherence term $\mathcal{L}_{\text{physics}}$
through the ALM framework, which introduces a Lagrange multiplier $\lambda_p$ and a quadratic penalty term with coefficient $\mu_p$.
Additionally, a separate energy conservation term $\mathcal{L}_{\text{energy}}$ is included with a weight $\lambda_E(k)$ that is dynamically
adjusted based on the curriculum stage.
The complete total loss is thus expressed as:

\begin{equation}
\mathcal{L}_{\text{total}} = \mathcal{L}_{\text{data}} + \lambda_p \mathcal{L}_{\text{physics}} + \frac{\mu_p}{2} \mathcal{L}_{\text{physics}}^2 + \lambda_E(k) \mathcal{L}_{\text{energy}}
\end{equation}


\textbf{ALM Initialization:}

The initial values for the ALM parameters are as follows:

\begin{itemize}
    \item $\lambda_p = 0.0$ \hspace{1em} (Initial Lagrange multiplier)
    \item $\mu_p = 0.01$ \hspace{1em} (Initial penalty coefficient)
    \item $v_p = 0.0$ \hspace{1em} (Exponential moving average of $\mathcal{L}_{\text{physics}}^2$)
    \item $\beta = 0.99$ \hspace{1em} (Momentum for $v_p$)
    \item $\gamma = 0.01$ \hspace{1em} (Base learning rate for $\mu_p$ update)
    \item $\epsilon = 1 \times 10^{-8}$ \hspace{1em} (Numerical stability constant)
\end{itemize}

\subsection[Data Fidelity Term: Anchoring to Ground Truth]{Data Fidelity Term ($\mathcal{L}_{\text{data}}$): Anchoring to Ground Truth}

The purpose of the data fidelity term is to ensure that the network's predictions align with the high-precision trajectory data generated by a trusted numerical integrator. This term grounds the general, physically-plausible solution learned by the PINN to the specific trajectory defined by a given set of initial conditions. While Mean Squared Error (MSE) is a common choice, Mean Absolute Error (MAE) is often more robust to the occasional large deviations or outliers that can occur in chaotic simulations. For this reason, the data fidelity loss is defined using the L1 norm$^1$:



The data fidelity loss is defined using the L1 norm on the \textbf{scaled positions}:

\[
\mathcal{L}_{\text{data}} = \frac{1}{N_{\text{data}}} \sum_{k=1}^{N_{\text{data}}} ||r'_{\text{pred}}(t'_k) - r'_{\text{true}}(t'_k)||_1
\]

where:
\begin{itemize}
    \item $t'_k$ is the normalized time coordinate in $[-1, 1]$
    \item $r'_{\text{pred}}(t'_k)$ is the network's predicted \textbf{scaled position} vector (the direct output of the network)
    \item $r'_{\text{true}}(t'_k)$ is the ground truth position from the BRUTUS data, \textbf{scaled by the characteristic length $D$} (i.e., $r'_{\text{true}} = r_{\text{true}}/D$)
    \item $N_{\text{data}}$ is the number of data points in a training batch
    \item $|| \cdot ||_1$ denotes the L1 norm
\end{itemize}





\textbf{Batch Construction:}

The data points are sampled from the BRUTUS trajectory snapshots as follows:

\begin{lstlisting}[language=Python, caption={Data Batch Construction}]
def get_batch(batch_size, points_per_trajectory=32):
    """
    Sample batch_size trajectories from the dataset.
    For each trajectory, sample points_per_trajectory time points.

    Returns:
        t_k: Time points sampled from BRUTUS snapshots
        r_true: Corresponding ground-truth positions

    Shapes: (batch_size * points_per_trajectory, ...)
    """
    # Implementation details depend on dataset structure
    return t_k, r_true

# Total data points per batch
N_data = batch_size * points_per_trajectory  # e.g., 256 * 32 = 8192
\end{lstlisting}

The time points $t_k$ are sampled from the discrete BRUTUS trajectory snapshots available in the training dataset, representing specific moments where high-precision numerical integration data is available.
\authcomment{Comment: This above batch construction is not consistent with the CTL instruction plan. Look into this.}



\textbf{Key ALM Integration Note:} \\
The data term is no longer weighted by a fixed parameter. Instead, its influence is automatically balanced against the physics constraint through the ALM framework. The effective weight of the physics term becomes $(\lambda_p + \mu_p \mathcal{L}_{\text{physics}})$, which adapts during training:
\begin{itemize}
    \item When physics loss is high $\rightarrow$ ALM increases $\lambda_p$ $\rightarrow$ physics term gets more weight
    \item When physics loss is low $\rightarrow$ ALM stabilizes $\lambda_p$ $\rightarrow$ balanced optimization
    \item The adaptive $\mu_p$ prevents oscillations by adjusting based on the variance of physics loss
\end{itemize}


\subsection[Physics Residual Term: Enforcing Newtonian Dynamics]{Physics Residual Term ($\mathcal{L}_{\text{physics}}$): Enforcing Newtonian Dynamics}

This term is the heart of the "physics-informed" paradigm. It enforces the governing laws of physics—in this case, Newton's law of universal gravitation—across the entire solution domain, not just at the discrete points where training data is available. This is what endows the network with the ability to generalize and make physically consistent predictions at any point in time.

The physics residual, $f_i(t')$, for each body $\mathrm{i}$ is defined as the discrepancy between the acceleration predicted by the neural network (via automatic differentiation) and the acceleration dictated by the gravitational forces from all other bodies $\mathrm{j}$:
\[
f_i(t') = a'_{\theta}(t', i) - \sum_{j \neq i} m_j \frac{r'_{\theta}(t', j) - r'_{\theta}(t', i)}{(\|r'_{\theta}(t', j) - r'_{\theta}(t', i)\| + \epsilon_{\text{soft}})^3}
\]

where:
\begin{itemize}
    \item $t'$ : Normalized time in $[-1, 1]$
    \item $r'$ : Positions scaled by $D$ (characteristic length)
    \item $a'$ : Accelerations scaled by $D/T^2$
    \item $m_j$ : Masses (unchanged, as $G=1$ in the dimensionless system)
    \item $\epsilon_{\text{soft}} = 10^{-5}$ : Softening parameter to prevent division by zero during close encounters
\end{itemize}

The physics loss is then formulated as the mean squared norm of this residual vector, summed over all bodies and evaluated over a set of $N_{\text{col}}$ collocation points:
\[
\mathcal{L}_{\text{physics}} = \frac{1}{N_{\text{col}}} \sum_{k=1}^{N_{\text{col}}} \sum_{i=1}^{N} || f_i(t'_k) ||^2
\]

\textbf{Collocation Point Sampling:}

The collocation points are sampled according to the following strategy:

\begin{lstlisting}[language=Python, caption={Collocation Points Sampling}]
import numpy as np

def sample_collocation_points(t_min, t_max, n_col, strategy='uniform',
                            traj_data=None):
    """
    Sample collocation points for physics and energy loss evaluation.

    Args:
        t_min, t_max: Normalized time domain boundaries (e.g., [-1, 1])
        n_col: Number of collocation points (recommended: 2048-4096)
        strategy: 'uniform' (default) or 'adaptive'
        traj_data: Optional trajectory data for adaptive strategy

    Returns:
        t_col: Array of sampled time points
    """
    if strategy == 'uniform':
        t_col = np.random.uniform(t_min, t_max, size=(n_col,))

    elif strategy == 'adaptive':
        if traj_data is None:  # This line needs proper indentation
            raise ValueError("Adaptive strategy requires traj_data")
        # Bias sampling toward high-acceleration regions
        weights = np.abs(traj_data['accelerations']).max(axis=1)

        # Safety check for degenerate cases
        if weights.sum() < 1e-8:
            weights = np.ones_like(weights)

        weights /= weights.sum()  # Normalize
        t_col = np.random.choice(traj_data['times'],
                                size=n_col, p=weights)
    else:
        raise ValueError("Invalid strategy")

    return t_col
\end{lstlisting}

Minimizing this term forces the network to learn a function $\mathrm{r}_\theta(t')$ that is also a solution to the governing ordinary differential equations (ODEs).

\textbf{ALM Implementation Details:}

The influence of the physics term is controlled adaptively through the Lagrange multiplier $\lambda_p$ and the penalty coefficient $\mu_p$. After each optimization step (primal step), these parameters are updated in a dual step as follows:

\begin{enumerate}
    \item \textbf{Update the Lagrange Multiplier ($\lambda_p$):} The multiplier is updated via gradient ascent, which naturally increases the penalty when the physics residual is large.

    \[
    \lambda_p \leftarrow \lambda_p + \mu_p \mathcal{L}_{\text{physics}}
    \]


    \item \textbf{Update the Adaptive Penalty ($\mu_p$):} The penalty coefficient is updated using a method inspired by the RMSprop optimizer to ensure stability.
    \begin{itemize}
        \item First, update a moving average of the squared physics loss, $v_p$:

        \[
        v_p \leftarrow \beta v_p + (1 - \beta) \mathcal{L}_{\text{physics}}^2
        \]

        \item Next, update the penalty parameter $\mu_p$. Based on the RMSprop analogy, the robust update rule is:

        \[
        \mu_p \leftarrow \frac{\gamma}{\sqrt{v_p + \epsilon}}
        \]
    \end{itemize}

    \item \textbf{Gradient Handling:} \\
    -- $\lambda_p$ must be treated as a non-trainable scalar updated via code (not autograd) \\
    -- $\mu_p$ is a dynamic parameter updated algorithmically (excluded from backpropagation)
\end{enumerate}

\begin{lstlisting}[language=Python, caption={ALM Parameter Updates}]
# In the training loop, after optimizer.step():
with torch.no_grad():
    # 1. Update Lagrange multiplier (gradient ascent)
    lambda_p += mu_p * physics_loss.item()

    # 2. Update moving average of squared physics loss
    v_p = beta * v_p + (1 - beta) * (physics_loss.item() ** 2)

    # 3. Update adaptive penalty (RMSprop-inspired)
    mu_p = gamma / (np.sqrt(v_p + epsilon))  # Note: using np.sqrt

    # 4. Clamp mu_p to prevent runaway values
    mu_p = np.clip(mu_p, 1e-4, 100)  # Changed from torch.clamp to np.clip

# Note: lambda_p, mu_p, v_p are NOT part of model.parameters()
# They are maintained as separate scalars with requires_grad=False
\end{lstlisting}

\subsection[HNN-Style Energy Regularizer: Enforcing a Global Conservation Law]{HNN-Style Energy Regularizer ($\mathcal{L}_{\text{energy}}$): Enforcing a Global Conservation Law}

While the physics residual enforces the local, instantaneous laws of motion, it does not explicitly guarantee the conservation of global quantities over long time horizons. For an isolated mechanical system like the $\mathrm{N}$-body problem, the total energy is one such conserved quantity. The energy loss term is designed to combat the long-term energy drift that can plague standard PINN models by explicitly penalizing deviations from this fundamental conservation law.$^1$

The total energy of the $\mathrm{N}$-body system, its Hamiltonian $H_{\text{total}}$, is the sum of the total kinetic energy $(T)$ and the total potential energy $(V)$:

\[
H_{\text{total}}(r', v') = T + V = \sum_{i=1}^N \frac{1}{2} m_i ||v'_i||^2 - \sum_{i<j} \frac{m_i m_j}{||r'_i - r'_j|| + \epsilon_{\text{soft}}}
\]


This represents the total energy in the non-dimensional unit system. The positions ($r'$) and velocities ($v'$) are the scaled outputs from the network.
The initial energy, $H_0$, must also be computed in these same dimensionless units for the loss term to be consistent.\\

The energy loss term penalizes any deviation of the predicted total energy from its initial total energy, $H_0$, which is a constant calculated from the known initial conditions. The loss is computed over the same set of $N_{\text{col}}$ collocation points as the physics loss:

\[
\mathcal{L}_{\text{energy}} = \frac{1}{N_{\text{col}}} \sum_{k=1}^{N_{\text{col}}} L_\delta(\Delta H_k)
\]

where:

\begin{enumerate}
    \item $\Delta H_k = H_{\text{total}}(r_\theta(t'_k), v_\theta(t'_k)) - H_0$

    \item $H_0$ is the precomputed initial energy for each trajectory (See the code below for the exact calculation.)

    \item $L_\delta$ is the Huber loss with $\delta = 1.0$:
    \[
    L_\delta(a) =
    \begin{cases}
        \frac{1}{2} a^2 & \text{for } |a| \leq \delta \\
        \delta (|a| - \frac{1}{2} \delta) & \text{for } |a| > \delta
    \end{cases}
    \]
\end{enumerate}

\textbf{Initial Energy Calculation:}

The initial energy $H_0$ is precomputed during dataset preprocessing:

\begin{lstlisting}[language=Python, caption={Initial Energy Calculation in Dataset Preprocessing}]
class TrajectoryDataset:
def __init__(self, trajectory_data):
    # ... other initialization ...

    # Precompute H_0 for each trajectory
    self.H_0_values = {}
    for traj_id, traj in trajectory_data.items():
        self.H_0_values[traj_id] = self.compute_initial_energy(
            r0_prime=traj['r0_scaled'],  # Already scaled by D
            v0_prime=traj['v0_scaled'],  # Already scaled by D/T
            masses=traj['masses']
        )

def compute_initial_energy(self, r0_prime, v0_prime, masses):
    """Compute H_0 in dimensionless units (G=1)"""
    H_0 = 0.0
    epsilon_soft = 1e-5  # Same softening parameter

    # Fixed for 3-body problem
    N_bodies = 3

    # 1. KINETIC ENERGY (dimensionless)
    for i in range(N_bodies):
        H_0 += 0.5 * masses[i] * np.sum(v0_prime[i]**2)

    # 2. POTENTIAL ENERGY (dimensionless)
    for i in range(N_bodies):
        for j in range(i+1, N_bodies):
            r_ij = np.linalg.norm(r0_prime[i] - r0_prime[j])
            H_0 -= masses[i] * masses[j] / (r_ij + epsilon_soft)

    return H_0
\end{lstlisting}

\textbf{Energy Regularization Schedule:}

The weight of the energy loss term is scheduled according to the curriculum stage:

\[
\lambda_E(k) = 10^{k-3} \quad \text{for } k \in \{0, 1, \ldots, 6\}
\]

This results in the following stage-specific values:
\begin{itemize}
    \item Stage 0: $\lambda_E = 10^{-3} = 0.001$
    \item Stage 1: $\lambda_E = 10^{-2} = 0.01$
    \item Stage 2: $\lambda_E = 10^{-1} = 0.1$
    \item Stage 3: $\lambda_E = 10^{0} = 1$
    \item Stage 4: $\lambda_E = 10^{1} = 10$
    \item Stage 5: $\lambda_E = 10^{2} = 100$
    \item Stage 6: $\lambda_E = 10^{3} = 1000$
\end{itemize}

\textbf{Key Implementation Notes:}
\begin{itemize}
    \item \textbf{Unit System Consistency:} All quantities (positions $r'$, velocities $v'$, and initial energy $H_0$) must be expressed in the same non-dimensional ($G = 1$) system of units.
    \item \textbf{Softening Parameter:} The same value $\epsilon_{\text{soft}} = 10^{-5}$ is used consistently in both the physics residual and energy calculations to prevent numerical instabilities during close encounters.
    \item \textbf{Collocation Points:} The same set of $N_{\text{col}}$ points is used for both physics and energy loss evaluation to ensure computational efficiency.
\end{itemize}
